\documentclass[]{article}

\usepackage[utf8]{inputenc}
\usepackage[T1]{fontenc}
\usepackage{lmodern}
\usepackage{textgreek}
\usepackage{newunicodechar}
\newunicodechar{⁰}{\ensuremath{^0}}
\newunicodechar{²}{\ensuremath{^2}}
\newunicodechar{³}{\ensuremath{^3}}
\newunicodechar{₁}{\ensuremath{_1}}
\newunicodechar{₂}{\ensuremath{_2}}
\newunicodechar{₃}{\ensuremath{_3}}
\newunicodechar{×}{\ensuremath{\times}}
\newunicodechar{∇}{\ensuremath{\nabla}}
\newunicodechar{≈}{\ensuremath{\approx}}
\newunicodechar{≡}{\ensuremath{\equiv}}
\newunicodechar{≤}{\ensuremath{\leq}}
\newunicodechar{≥}{\ensuremath{\geq}}
\newunicodechar{→}{\ensuremath{\rightarrow}}
\newunicodechar{←}{\ensuremath{\leftarrow}}
\newunicodechar{↑}{\ensuremath{\uparrow}}
\newunicodechar{↓}{\ensuremath{\downarrow}}
\newunicodechar{↔}{\ensuremath{\leftrightarrow}}
\newunicodechar{—}{\textemdash}
\newunicodechar{–}{\textendash}
\newunicodechar{−}{\ensuremath{-}}
\newunicodechar{√}{\ensuremath{\surd}}
\newunicodechar{∼}{\ensuremath{\sim}}
\newunicodechar{≲}{\ensuremath{\lesssim}}
\newunicodechar{≳}{\ensuremath{\gtrsim}}
\newunicodechar{⟨}{\ensuremath{\langle}}
\newunicodechar{⟩}{\ensuremath{\rangle}}
\newunicodechar{₀}{\ensuremath{_0}}
\newunicodechar{₄}{\ensuremath{_4}}
\newunicodechar{₅}{\ensuremath{_5}}
\newunicodechar{₆}{\ensuremath{_6}}
\newunicodechar{₇}{\ensuremath{_7}}
\newunicodechar{₈}{\ensuremath{_8}}
\newunicodechar{₉}{\ensuremath{_9}}
\newunicodechar{⁴}{\ensuremath{^4}}
\newunicodechar{α}{\ensuremath{\alpha}}
\newunicodechar{β}{\ensuremath{\beta}}
\newunicodechar{γ}{\ensuremath{\gamma}}
\newunicodechar{δ}{\ensuremath{\delta}}
\newunicodechar{ε}{\ensuremath{\epsilon}}
\newunicodechar{θ}{\ensuremath{\theta}}
\newunicodechar{λ}{\ensuremath{\lambda}}
\newunicodechar{μ}{\ensuremath{\mu}}
\newunicodechar{ν}{\ensuremath{\nu}}
\newunicodechar{ξ}{\ensuremath{\xi}}
\newunicodechar{π}{\ensuremath{\pi}}
\newunicodechar{ρ}{\ensuremath{\rho}}
\newunicodechar{σ}{\ensuremath{\sigma}}
\newunicodechar{τ}{\ensuremath{\tau}}
\newunicodechar{φ}{\ensuremath{\phi}}
\newunicodechar{ϕ}{\ensuremath{\phi}}
\newunicodechar{χ}{\ensuremath{\chi}}
\newunicodechar{ψ}{\ensuremath{\psi}}
\newunicodechar{ω}{\ensuremath{\omega}}
\newunicodechar{Δ}{\ensuremath{\Delta}}
\newunicodechar{Λ}{\ensuremath{\Lambda}}
\newunicodechar{Σ}{\ensuremath{\Sigma}}
\newunicodechar{Φ}{\ensuremath{\Phi}}
\newunicodechar{Ω}{\ensuremath{\Omega}}
\newunicodechar{“}{``}
\newunicodechar{”}{''}
\newunicodechar{‘}{`}
\newunicodechar{’}{'}
\usepackage{amsmath}
\usepackage{amssymb}
\usepackage{multirow}
\usepackage{natbib}
\usepackage{graphicx}
\usepackage{tabularx}
\usepackage{adjustbox}
\usepackage{listings}
\usepackage{hyperref}
\usepackage{xcolor}
\lstset{
  basicstyle=\ttfamily\footnotesize,
  backgroundcolor=\color{black!4},
  frame=none,
  breaklines=true,
  breakatwhitespace=true,
  showstringspaces=false,
  columns=fullflexible,
  xleftmargin=1em,
  xrightmargin=1em,
  aboveskip=0.6em,
  belowskip=0.6em,
  keywordstyle=\color{blue!60!black},
  commentstyle=\itshape\color{black!60},
  stringstyle=\color{green!50!black},
}
\usepackage[margin=1in]{geometry}

\begin{document}

In this work, we introduced a graph-latent residual decomposition framework to rigorously isolate and quantify non-additive electronic effects from the extensive scaling inherent in molecular internal energy ($U_0$). By first establishing a classical group-contribution baseline using Ridge regression, we successfully decoupled size-dependent energetic scaling. The remaining residuals, representing purely non-linear quantum-mechanical deviations, were then modeled using a Graph Neural Network (GNN). Our results demonstrate that explicit graph-based message passing is highly effective at capturing these topological nuances. The GNN achieved a test $R^2$ of 0.9985 on the standardized residuals, substantially outperforming a descriptor-based Random Forest baseline and confirming that molecular topology is essential for modeling complex electronic physics.Beyond predictive accuracy, this framework provides deep interpretability into the learned representations of non-additivity. Analysis of the GNN's 128-dimensional latent space revealed that the network intrinsically learned to encode fundamental intensive properties, such as the HOMO-LUMO gap and HOMO energy, despite receiving no direct supervision on these targets. Dimensionality reduction further illustrated a highly structured embedding landscape where distinct property gradients align with the modeled energetic deviations, proving that the network captures genuine quantum-mechanical phenomena rather than arbitrary statistical variance.Furthermore, by applying Integrated Gradients, we successfully mapped these non-additive energetic deviations back to specific structural motifs \citep{sundararajan2017axiomaticattributiondeepnetworks,polat2025xchemagentsagenticaiexplainable,ji2026molledgeradditivegraphneural}. Statistical validation confirmed that features such as aromatic rings and the integration of highly electronegative heteroatoms (nitrogen and fluorine) are significant drivers of non-additivity \citep{lundstrom2022rigorousstudyintegratedgradients,wang2026explainablemolecularpropertyprediction}. These motifs correspond directly to physically observable shifts, notably compressed HOMO-LUMO gaps and increased dipole moments, which are characteristic of extended $\pi$-conjugation and complex electron localization \citep{chen2026interpretablemachinelearningquantuminformed,mick2026datafusioncontrastivealignment,bigting2026martinimapperautomatedfragmentbased}.Ultimately, this approach bridges the gap between classical additivity models and advanced deep learning techniques. By explicitly separating extensive scaling from underlying non-linear electronic physics, our framework not only improves predictive performance but also provides a chemically intuitive mapping of how specific structural motifs deviate from classical group-contribution theory. This methodology offers a robust and interpretable tool for future computational chemistry applications, enabling more targeted molecular design and a deeper understanding of structure-property relationships without the computational overhead of 3D coordinate generation.

\end{document}
                